% !TEX root = ../main.tex
\section{Introduction}
\label{sec:intro}

\begin{table}[t!]
\centering
\small
\setlength{\tabcolsep}{4pt}
\begin{tabular}{lccccc}
\toprule
signal & vs.\ human $r$ & Spearman $\rho$ & within-correct $r$ & ranking (hard/sat.) & item share \\
\midrule
\multicolumn{6}{l}{\emph{\footnotesize the free baseline}} \\
\textbf{free answer-correct bit} & $\mathbf{+0.92}$ & --- & (constant) & --- & --- \\
\midrule
\multicolumn{6}{l}{\emph{\footnotesize verification signals (use the reference only as an endpoint)}} \\
completion curve ($k{=}0.5$) & $+0.76$ & $+0.48$ & $+0.02$ & --- & --- \\
completion curve ($k{=}0.75$) & $+0.78$ & $+0.46$ & $+0.19$ & --- & --- \\
completion-curve mean & $+0.79$ & $+0.41$ & $+0.07$ & --- & --- \\
\midrule
\multicolumn{6}{l}{\emph{\footnotesize matching metrics (score resemblance to the reference text)}} \\
\tfa{} (argmax) & $-0.16$ & $-0.12$ & $+0.21$ & $+0.60$ / $-0.25$ & $83$--$92\%$ \\
gold-token logprob & $-0.27$ & $-0.18$ & $+0.15$ & $+0.43$ / $-0.23$ & $83$--$91\%$ \\
TF-IDF cosine & $-0.08$ & $-0.09$ & $+0.16$ & $+0.72$ / $+0.12$ & $69\%$ \\
sentence embeddings & $\mathbf{+0.14}$ & $+0.13$ & $-0.01$ & $+0.87$ / $+0.29$ & $\sim$$70\%$ \\
unigram-F1 & $\mathbf{-0.33}$ & $-0.29$ & $+0.06$ & $-0.49$ / $+0.18$ & $45$--$55\%$ \\
\midrule
\multicolumn{6}{l}{\emph{\footnotesize LLM judges (expensive comparison points)}} \\
LLM judge (referenced) & $+0.89$ & $+0.43$ & $-0.14$ & --- & --- \\
LLM critic (GPT) & $+0.91$ & $+0.61$ & $+0.17$ & --- & --- \\
LLM critic (Claude) & $+0.88$ & $+0.54$ & $+0.03$ & --- & --- \\
\bottomrule
\end{tabular}
\caption{\textbf{No cheap signal grades beyond the free correctness bit: every matching metric lands in $[-0.33, +0.14]$ against blind human soundness, while the free answer-correct bit tops the entire field at $+0.92$.}
Each row is one evaluation signal scored against blind human soundness on the anchor set of 60 human-graded solutions: Pearson and Spearman correlations, the within-correct correlation (restricted to solutions whose final answer is correct, the margin a grader must win), the model-ranking Spearman on hard versus saturated benchmarks, and the share of per-solution variance owned by the item (the reference).
The completion curve truncates a solution, lets a fixed completer finish it, and scores whether the completion reaches the right answer; it adds no measurable grading signal beyond the free bit (Section~\ref{sec:verify}).
Point estimates only; confidence intervals, subset sizes, and per-row caveats are in the twin table, Table~\ref{tab:citwin} (Appendix~\ref{app:citwin}).}
\label{tab:classfull}
\end{table}

% P1: worked example
A benchmark ships a worked reference solution for every problem, and the cheapest grader scores a model's solution by its similarity to that reference.
A student who writes a beautiful wrong derivation matches the reference well; a student who solves the problem an unusual correct way matches it badly.
We measure exactly what similarity rewards, across seventeen models from four families on eight math and code benchmarks, anchored to blind human grades and to executed tests.
The answer is mostly ``how predictable the reference is'': most of a matching score is a property of the reference text, not of the solution being graded.
The only reliable cheap check turns out to be the one everyone already has, comparing final answers.

% P2: contributions
Concretely, this paper contributes:
\begin{itemize}
\item \textbf{A grading null.} Five matching metrics spanning the space, from teacher-forced accuracy to sentence embeddings, all land in $[-0.33, +0.14]$ against blind human grades of reasoning soundness, and the null replicates on code against executed tests; the null is structural for the two prediction-fit metrics, which never read the generated solution, and empirical for the three output-similarity metrics that do (Section~\ref{sec:classnull}, Table~\ref{tab:classfull}).
\item \textbf{A mechanism.} Most of a matching score's per-solution variance is an item effect that replicates across model families, so the score measures how predictable the reference is rather than grading the solution in front of it (Section~\ref{sec:mechanism}, Figure~\ref{fig:mechanism}).
\item \textbf{A ranking boundary.} Matching can still rank models, but only in aggregate, on hard benchmarks, with worked references; replacing each reference with its bare answer line collapses the ranking on the same items, and the registered length-matched control (P-L1) has since run and landed as predicted, so the collapse tracks reference informativeness rather than token count (Section~\ref{sec:ranking}, Table~\ref{tab:map}).
\item \textbf{A free baseline that nothing beats.} The machine-checked answer-correct bit alone correlates $+0.92$ with human soundness, a level expected from construct overlap and treated as the baseline to beat rather than as a finding, and on the anchor set no signal we tested, matching or verification, adds measurable grading signal beyond it (Section~\ref{sec:verify}).
\item \textbf{A registered record.} Six predictions were registered before their data were generated, and four of six landed as called; the complete ledger, including one sign flip and one falsification, is Table~\ref{tab:ledger} (Section~\ref{sec:slate}).
\end{itemize}

% P3: provenance and framing
The sharpest matching metric is teacher-forced accuracy (\tfa): feed the model the problem and the gold solution one token at a time, and score the fraction of gold tokens the model would have predicted itself.
\tfa{} was introduced as a judge-free evaluation signal by \textsc{Putnam-AXIOM}~\cite{gulati2025putnam}; we do not introduce it here, we characterize it, mapping where it works, where it saturates, and where it fails.
The appeal of matching is real: it is deterministic, costs one forward pass, involves no judge prompt whose wording could shift the score, and anyone can audit it, on benchmarks from grade-school math to program synthesis~\cite{cobbe2021gsm8k,hendrycks2021math,chen2021codex}.
Prior work organizes this space along a single axis of cost versus quality, string and embedding metrics at the cheap end~\cite{yuan2021bartscore} and LLM judges at the expensive end~\cite{zheng2023judging}, with the practical question being how far down the cost axis one can afford to go.
That framing hides an assumption: that similarity degrades gracefully, so a cheaper metric measures the same thing with more noise.
The assumption fails, because the operative axis is not cost but mechanism: signals that score resemblance to the reference text (matching) behave differently in kind from signals that check whether the solution reaches a checkable endpoint (verification), and in every configuration we tested, matching met no grading standard at any price point we tried.

% P4: anchoring
Everything is anchored to blind human grades on two deliberately separated scales, soundness of the reasoning and quality of the writing, plus executed tests on code, where no human is needed.
The grades carry a measured reliability: a hypothesis-blind rater outside the author list regraded half the anchor items under the same rubric and agreed at $0.83$ on soundness (Section~\ref{sec:setup}, Appendix~\ref{app:ceiling}).
If a metric is validated only by its agreement with an LLM judge, that is reliability, not validity; if the judge scores surface style, a metric tracking it may be succeeding at the wrong thing.
We test five matching metrics, targeted repairs of each, and three judge-free verification signals: executed tests, a completion curve that truncates a solution and checks whether a fixed completer still reaches the right answer, and an answer-swap check.
Table~\ref{tab:classfull} shows the full field of signals on one page.

% P5: the four questions
The results answer four questions.
\textbf{RQ1}: can matching grade individual solutions? No: all five metrics land in the band above, structurally so for the prediction-fit metrics, which never read the solution, and empirically so for the output-similarity metrics; a frontier LLM judge reaches $+0.89$ on identical items, and on code the prediction-fit metrics are indistinguishable from zero against executed tests (Section~\ref{sec:classnull}).
\textbf{RQ2}: why not? Because $83$--$92\%$ of \tfa's per-solution variance beyond model level (a share stable under item resampling and single-model deletion) is an item effect that replicates across families, and pooled across model--item pairs \tfa{} is largely reference perplexity; reference predictability and problem difficulty are entangled inside that item effect, so the score meters a text property either way, and does not grade the solution in front of it (Section~\ref{sec:mechanism}).
\textbf{RQ3}: when can matching still rank models? On hard benchmarks with information-dense references, where every family size ladder is positive, and not on saturated ones, where five of twelve ladders invert; every targeted repair we or our reviewers proposed left the per-solution signal at zero or made it worse (Sections~\ref{sec:ranking} and~\ref{sec:slate}).
\textbf{RQ4}: what does grade? The free answer-correct bit, and on the anchor set nothing we tested beats it: the completion curve comes closest yet adds no measurable signal beyond the bit, its registered completer-strength prediction was falsified by a flat curve across a twenty-fold completer-size range, and a cheap beyond-correctness grader remains an open problem (Section~\ref{sec:verify}).

% P6: implications
These results replace a folk cost--quality tradeoff with a negative map, a mechanism, and an advance-registered test of both.
In the entire space we tested, no cheap signal detectably grades beyond correctness, with within-correct effects above roughly $0.4$ excluded by equivalence testing, and correctness itself is free where references exist.
For benchmark builders, a reference's value is as an endpoint for a checker, or as context for a reader, never as a similarity target.
For leaderboards, matching is legitimate exactly when references are worth predicting, and reference perplexity can be measured in advance to check that condition.
